% \tracingmacros DRAWS THE MACRO BEING EXPANDED, AND ITS ARGUMENTS.
%
% Before a macro is expanded the reference draws a blank line, the
% macro's name, its parameter text, `->' and its body -- the same text
% \show gives after the `macro:' it puts in front:
%
%   (blank)
%   \d #1\d ->#1#1
%
% and then one line per argument as it is matched, with no blank line of
% its own:
%
%   #1<-##
%
% A parameter character inside an argument is shown doubled, as it is
% everywhere a token list is shown. The name carries the space a control
% sequence takes when something follows it, so a control word and a
% single letter both read `\A ->B' rather than `\A->B'.
%
% One draws at \tracingmacros=1 as well as at 2; the difference between
% the two is the token PARAMETERS -- \output, \everyvbox and their kin --
% which are drawn only at 2 and which HSTeX does not draw yet.
\catcode`\{=1 \catcode`\}=2 \catcode`\#=6
\tracingonline=1
\def\A{B}
\def\d#1\d{#1#1}
\long\def\L#1{[#1]}
\outer\def\O{o}
\def\two#1#2{#2#1}
\message{[A a macro with no parameters]}
\tracingmacros=2
\message{\A}
\tracingmacros=0
\message{[B a delimited parameter holding a hash]}
\tracingmacros=2
\message{\d#\d}
\tracingmacros=0
\message{[C two parameters]}
\tracingmacros=2
\message{\two XY}
\tracingmacros=0
\message{[D a long macro with a braced argument]}
\tracingmacros=2
\message{\L{a b}}
\tracingmacros=0
\message{[E at tracingmacros=1]}
\tracingmacros=1
\message{\two XY}
\tracingmacros=0
\message{[done]}
\end
